\chapter{一致空间.}

涉及一致空间的主要概念和命题是从实变量理论中一点点脱胎出来的，直到晚近才成为系统研究的对象。柯西为了严格地建立级数理论（参见第 153 页），把一条他似乎视为显然的原则当作出发点：序列 \((a_{n})\) 收敛的一个必要充分条件是，\(|a_{n+p} - a_{n}|\) 只要当 \(n\) 充分大就可任意小（例如见（{[}56 a{]}，(2)，第 VII 卷，第 267 页））。他与波尔查诺{[}27 c{]}无疑都属于最早明确陈述这一原则并认识到其重要性的人：因此，满足上述条件的实数序列被称为“柯西序列”，而且引申开来，度量空间中这样的点列 \((x_{n})\) 也获此名：\(x_{n+p}\) 与 \(x_{n}\) 之间的距离只要当 \(n\) 充分大就可任意小；最终，一致空间中柯西序列的推广便得到了“柯西滤子”之名。

后来，当实数的直观概念不再令人满意、人们为了给分析以坚实的基础而设法从有理数出发定义实数时，在 19 世纪下半叶提出的各种定义中，恰恰是柯西原则提供了最富成效的一种；这就是康托尔的定义（{[}47{]}，第 93–96 页）（海涅{[}154 b{]}等人也从康托尔的思想出发发展了它，梅雷则独立地作了发展），按照这个定义，让每个有理数的柯西序列（用康托尔的术语说，即“基本序列”）对应于一个实数；两个有理数柯西序列 \((a_{n})\) 与 \((b_{n})\) 对应于同一个实数，当且仅当 \(|a_{n} - b_{n}|\) 趋于零。这里的本质思想是，从某种观点看，有理数集 \(\mathbf{Q}\) 是“不完备的”，而实数集是由 \(\mathbf{Q}\) 经“完备化”而导出的那个“完备”集。

另一方面，海涅在主要受魏尔斯特拉斯和康托尔的思想启发的工作中，第一个定义了一元或多元实变量数值函数的一致连续性{[}154 a{]}，并证明了在 R 的一个闭有界区间上连续的每个数值函数都在其上一致连续{[}154 b{]}：这就是“海涅定理”。这一结果联系着 R 中闭有界区间的紧性（“博雷尔–勒贝格定理”，参见第 141 页），而且海涅对其定理所给的证明略作修改也可用来证明博雷尔–勒贝格定理（在好几位作者看来，这就足以成为把这个定理命名为“海涅–博雷尔定理”的理由）。

这些思想向更一般空间的推广，是在人们先研究一些特殊情形、然后一般地研究度量空间时完成的；在度量空间中给定了距离（点对的数值函数，满足某些公理），它同时定义了一个拓扑和一种一致结构。弗雷歇第一个陈述了这些空间的一般定义，认识到了柯西原则的重要性{[}115 a{]}，并且还就一切度量空间引入了预紧（或“全有界”{[}115 a 与 b{]}）空间的概念。豪斯多夫在其《集合论》（Mengenlehre）{[}152 a 与 b{]}中发展了度量空间理论的许多内容，他特别认识到，可以把上面考察过的康托尔构造应用于这些空间，从而对每个非“完备”的度量空间（也就是说柯西原则不成立的度量空间）导出一个“完备”的度量空间。

度量空间是一种特殊类型的“一致空间”；一致空间直到最近才由 A. 韦伊{[}330 b{]}以一般的方式加以定义。在此之前，人们只知道在处理度量空间时如何使用涉及“一致结构”的概念和结果：这就解释了在现代拓扑学的许多工作中，度量空间或可度量化空间（特别是紧的可度量化空间）为何在一些距离并无真正用处的问题里扮演着重要角色。一旦陈述了一致空间的定义，把例如豪斯多夫所陈述的度量空间理论的几乎全部内容推广到这些空间，便没有任何困难（尤其是当滤子的概念也可供使用时）（而且可以用同样的方式，把亚历山德罗夫–霍普夫的《拓扑学》{[}4{]}中就紧度量空间所陈述的结果，例如推广到一切紧空间）。特别地，一致空间的完备化定理无非是康托尔关于实数的构造的移植，不包含任何本质的修改。

